\documentclass[11pt,a4paper]{article}
\usepackage[T1]{fontenc}
\usepackage{lmodern,amsmath,amssymb,amsthm,hyperref}
\usepackage[margin=26mm]{geometry}
\newtheorem{lemma}{Lemma}
\newtheorem{theorem}[lemma]{Theorem}
\newcommand{\E}{\mathbb E}
\newcommand{\Prob}{\mathbb P}
\title{Proof Audit of Finite Global Realization}
\author{}
\date{4 October 2026}
\begin{document}
\maketitle
\noindent\textbf{Scope and verdict.}
This note audits Lemma 2 and its stochastic consequence in the supplied
\texttt{causal\_spectrum\_complete(1).tex}. Within the stated finite,
full-history-indexed causal model, the argument is valid. No missing
asymptotic estimate or unproved realization assumption is needed for this
lemma. We expand its minimax step, adaptive correctness, and the treatment
of distributional atoms. This is an additional mathematical audit, not
independent external peer review or a proof-assistant certificate. It does
not audit all later asymptotic sections or establish priority.

\section{Objects and exact hypotheses}
Let a finite rooted tree have action/output children $vAy$. Each nonterminal
node has a nonempty finite action set, and every action has a nonempty
finite output set. Each node has a unique complete history. All paths end
at terminal nodes. Allowing history-dependent action sets or unequal terminal
depths only generalizes the manuscript's setting.
A nonnegative flow $r$ satisfies
\[
\sum_y r(vAy)=r(v)\quad\text{for every nonterminal }v\text{ and action }A.
\]
A response strategy chooses one output for each history/action pair. Its
compatible histories have indicator $\chi_d$, itself a flow of root mass
one. A deterministic policy chooses one action at each history; its
terminal set is $\mathcal L_\pi$. A strategy and policy have exactly one
common terminal history. Unreachable choices may always be completed
arbitrarily, so strategies are genuine total decoders.

\section{The minimax step}
For any nonnegative flow set
\[
b_r(v)=r(v)\quad(v\text{ terminal}),\qquad
b_r(v)=\min_A\max_y b_r(vAy)\quad(v\text{ nonterminal}).
\]
\begin{lemma}[Explicit tree minimax]
\[
\max_d\min_{\substack{v\text{ terminal}\\\chi_d(v)=1}}r(v)
=b_r(\varnothing)
=\min_\pi\max_{v\in\mathcal L_\pi}r(v).
\tag{1}
\]
This is also the largest $w$ for which $r-w\chi_d\ge0$ at every node for
some strategy $d$.
\end{lemma}
\begin{proof}
For the left side, a strategy specifies an independent output choice and
subtree strategy for every available action. Its worst compatible terminal
weight is the minimum over actions of the selected subtree minima.
Backward induction therefore gives $\min_A\max_y b_r(vAy)$.
For the right side, a policy chooses one action and specifies a continuation
policy separately after each possible output. Its largest terminal weight
is the maximum of those subtree maxima. Minimizing the independent
continuations gives the same recursion.

A signed flow is nonnegative at every node if it is nonnegative at every
terminal node: induction follows from any action's conservation equation.
Since $r-w\chi_d$ is a signed flow, the terminal conditions are sufficient
as well as necessary. Finally, if $r(v)>0$, each action has a positive
child; induction gives $b_r(v)>0$. Zero root flow forces all nodes to have
zero weight.
\end{proof}
In particular, choosing a minimizing action at each node gives a policy
whose every terminal residual is at most $b_r(\varnothing)$. This statement
concerns the residual leaf weights, not their normalization as a new
probability model. No normalization is needed.

\section{Extraction and simultaneous correctness}
Start at a flow $p$ of root mass one. At each nonzero residual extract a
strategy of weight $w=b_r(\varnothing)>0$ as in (1). The residual remains
a nonnegative flow. A compatible terminal node attains the minimum $w$,
so at least one previously positive terminal weight becomes zero. Zero
weights can never increase. Consequently at most $L_+$ extractions occur,
where $L_+$ is the number of initially positive terminal histories.
The procedure ends with zero residual at every node and gives
\[
p(v)=\sum_jw_j\chi_{d_j}(v),\qquad \sum_jw_j=1.
\tag{2}
\]
The bottleneck is coordinatewise nondecreasing in residual weights, so the
extracted weights are nonincreasing. Real-valued or irrational kernels
cause no difficulty: finiteness follows from killed leaves, not a rational
lattice or a numerical tolerance.

Choose one $d_j$ with probability $w_j$. Equation (2) realizes every history
cylinder simultaneously. For a deterministic external policy, on its
reachable histories the probability of history $v$ is $p(v)$, and the
probability of its child $vAy$ is $p(vAy)$. Their ratio is the specified
kernel whenever $p(v)>0$. A randomized policy independent of the seed can
be conditioned on its private randomness; each conditioned policy has the
same correctness property. This proves adaptive correctness. At zero-mass
histories conditional laws impose no positive-probability constraint.

This step uses distinct coordinates for distinct complete histories. It
does not prove a corresponding result when decoders at different histories
are required to share the same time/action response coordinate.

\section{The decisive soft-tail inequality}
For each policy let $I_\pi$ be the surprisal of its random terminal history.
Zero-probability leaves are omitted when taking logarithms. Since there
are finitely many policies,
\[
F_*(t)=\min_\pi\Prob(I_\pi\le t)
\]
is a finite-support CDF of a nonnegative random variable $Z$. Thus every
$I_\pi$ is stochastically dominated by $Z$.

Fix $0<\delta\le1$. Stop extraction immediately before the first weight
at most $\delta$, or at termination if no such weight exists. Nonincreasing
weights show that the remaining root mass is
\[
r(\varnothing)=\sum_{j:w_j\le\delta}w_j.
\]
If it is nonzero, (1) supplies a policy with $r(v)\le\delta$ at every
terminal history. Moreover $r(v)\le p(v)$. Policy conservation and stochastic
dominance now give
\begin{align*}
\sum_{j:w_j\le\delta}w_j
&=\sum_{v\in\mathcal L_{\pi_r}}r(v)\\
&\le\sum_{v\in\mathcal L_{\pi_r}}\min\{p(v),\delta\}\\
&\le\max_\pi\E\min\{1,\delta2^{I_\pi}\}\\
&\le\E\min\{1,\delta2^Z\}.
\end{align*}
The zero-residual case is immediate. The function of information in the
last line is bounded and increasing, precisely the direction needed for
the stochastic comparison. Zero leaf weights contribute zero to the
clipped sum and require no convention involving $0\cdot\infty$.

\section{Stochastic sandwich and entropy}
Write $J_G=-\log_2 w_{S_G}$, and take an auxiliary
$E\sim\mathrm{Exp}(\ln2)$ independent of $Z$. For $t\ge0$,
\[
\Prob(J_G\ge t)
\le\E\min\{1,2^{Z-t}\}
=\Prob(Z+E\ge t).
\tag{3}
\]
The equality holds also when $t$ equals an atom of $Z$: conditional on
$Z\ge t$, the exponential tail is one. For $t<0$, all variables are
nonnegative. To pass from the inclusive tail inequality to CDFs, take
thresholds $t+1/k$ and let $k\to\infty$. This gives
$\Prob(J_G>t)\le\Prob(Z+E>t)$, and hence
$J_G\preceq_{\rm st}Z+E$ without dropping or adding any atom.

For any exact finite seed $S$ and deterministic policy $\pi$, a seed atom
lies inside the resulting transcript atom. Consequently, on their induced
joint space, $J_S\ge I_\pi$. Therefore
\[
F_{J_S}\le F_{I_\pi}\quad\text{for all }\pi,
\qquad F_{J_S}\le F_*,
\]
which is $Z\preceq_{\rm st}J_S$. We have proved
\begin{theorem}
For every such finite controlled tree there exists an exact seed with at
most $L_+$ atoms satisfying
\[
\boxed{Z\preceq_{\rm st}J_G\preceq_{\rm st}Z+E.}
\]
Every exact seed satisfies the first inequality. In particular,
\[
\boxed{\E Z\le C\le H(S_G)\le\E Z+\log_2e.}
\]
\end{theorem}
\begin{proof}
Only the expectation statement remains. Integrating nonnegative-variable
tails in the stochastic sandwich gives it, since $\E E=1/\ln2$.
\end{proof}
Independence of $E$ and $Z$ is an auxiliary comparison construction, not
an assertion that the greedy information equals $Z$ plus independent
operational noise. Equal action entropies, a CLT, FGL, and a growing horizon
are not used anywhere in this argument. For a depth-zero tree the unique
seed has weight one, $Z=J_G=0$, and the statements hold directly.

\section{Independent finite checks}
\texttt{audit\_incidence.py} enumerates strategy and policy sets as subsets
of terminal histories. It does not call or reproduce the manuscript's
bottom-up bottleneck routine. At each extraction it compares directly
\[
\max_d\min_{v\in d}r(v)
\quad\text{and}\quad
\min_\pi\max_{v\in\pi}r(v).
\]
It also checks unique strategy/policy intersection, all residual policy
masses, exact reconstruction, monotone extraction weights, seed support,
the spectrum converse, and both soft-tail inequalities using fractions.
Deterministic transitions and zero-weight branches are retained.

The run covers 165 cases, 769 residual minimax checks, and 2757 probability
threshold checks, with no floating-point arithmetic. The cases include
horizon zero, fair and nonuniform laws, and reproducibly generated
history-dependent laws. All passed. These are finite diagnostic checks;
the universal proof is the argument above. Run
\begin{verbatim}
python -S audit_incidence.py
\end{verbatim}
The JSON report explicitly records that no proof-assistant certificate is
provided. No larger-horizon testing is needed to justify the general proof.

\section{Editorial disposition}
No corrective mathematical change to Lemma 2 is required by this audit.
For readability, the original manuscript could display (1) explicitly,
state that unreachable decoder choices are completed arbitrarily, and
include the one-sided-limit sentence following (3). These make existing
steps explicit; they do not supply missing hypotheses or repair a failed
inequality. Questions of priority, sharpness of the additive constant,
and the equal-varentropy logarithmic lower bound remain separate.
\end{document}
